\section{Readable records and retained route coherence}
\label{sec:chvif-readable-records}

The following statements use the quantum state, channel, and measurement
premises already declared in Chapter VI. They make the relation between a
readable record and an unsorted interference experiment quantitative. The
record labels must refer to the same routes and the same ensemble as the
visibility; a reliable record of an unrelated label supplies no such bound.

\subsection{A same-state, worst-label theorem}
\label{sec:chvif-record-visibility-theorem}

Let $S$ be the two-route system and let $M$ contain the remaining physical
marker degrees of freedom. At the chosen observation time write an arbitrary
normalized joint state as
\begin{equation}
\begin{gathered}
\rho_{SM}=\sum_{j,k=0}^{1}|j\rangle\langle k|\otimes R_{jk},
 \\
 p=\operatorname{Tr}R_{00},\quad q=\operatorname{Tr}R_{11}=1-p.
\end{gathered}
\label{eq:chvif-record-block-state}
\end{equation}
For $p,q>0$, the conditional marker states are
$\sigma_M^0=R_{00}/p$ and $\sigma_M^1=R_{11}/q$. Their labels specify
the two diagonal route alternatives used by this operational test; this
notation does not select a unique outcome in a coherent joint state.

Declare a label-independent channel $\Lambda_t:M\longrightarrow D$ that
includes the storage and access operations available to the reading domain
at time $t$. It is completely positive and trace preserving, and acts on
the marker rather than on $S$. Let $F_0(t),F_1(t)$ be one physically
allowed decoder POVM on $D$, with $F_0+F_1=I_D$. The decoder may depend
on the declared time and controls, but receives no copy of the unknown
label. Define
\begin{equation}
\begin{gathered}
\tau_D^j(t)=\Lambda_t(\sigma_M^j),\\
 r_j(t)=\operatorname{Tr}[F_j(t)\tau_D^j(t)].
\end{gathered}
\label{eq:chvif-decoder-conditional-success}
\end{equation}
Storage and access can instead be specified jointly on an enlarged marker
containing a correlated environment, a transport carrier, and retained
classical records. No product-state or independent-noise assumption is
required. An environment correlated with the marker must be included in
that joint channel description when the proposed reduced channel would
otherwise depend on its correlations. Conditional states for the theorem
are always the actual conditional states of
Eq.~\eqref{eq:chvif-record-block-state} under that description.

Let $c=\langle0|\rho_S|1\rangle$. For a lossless two-route preparation
and the declared balanced, variable-phase analyser, with no
outcome-dependent filtering, the unsorted single-port visibility is
\begin{equation}
 V=2|c|.
 \label{eq:chvif-unsorted-record-visibility}
\end{equation}
The local, unselected channel $\Lambda_t$ does not change $\rho_S$.
Consequently a decoder can be analysed on $D$ while the visibility refers
to this same route marginal. A joint controlled inverse or feed-forward
operation on $S$ would define a different experiment and must be included
before applying the theorem to its resulting state.

\paragraph{Theorem.}
If, at the chosen time, both conditional decoder successes satisfy
\begin{equation}
\begin{gathered}
\min\{r_0(t),r_1(t)\}\geq1-\varepsilon,
 \\
 0\leq\varepsilon<\tfrac12,
\end{gathered}
\label{eq:chvif-worst-label-readability}
\end{equation}
then
\begin{equation}
 \boxed{\begin{gathered}
 V\leq4\sqrt{pq\,\varepsilon(1-\varepsilon)}
       \leq2\sqrt{\varepsilon(1-\varepsilon)} .\end{gathered}}
 \label{eq:chvif-readable-visibility-bound}
\end{equation}
For balanced priors the two bounds agree. For unequal priors the first
retains their additional restriction on the fringe. At $p=0$ or $q=0$,
positivity already gives $c=0$ and $V=0$; a conditional success statement
about an absent label must be supplied by a separately declared calibration
preparation if it is wanted.

\paragraph{Derivation.}
Apply the decoder to $(I_S\otimes\Lambda_t)(\rho_{SM})$ and retain its
two outcomes. The unnormalized route state for outcome $y$ is positive,
with diagonal entries $p a_y,q b_y$, where
\begin{equation}
\begin{gathered}
a_y=\operatorname{Tr}[F_y\tau_D^0],\\
 b_y=\operatorname{Tr}[F_y\tau_D^1].
\end{gathered}
\end{equation}
If its off-diagonal entry is $c_y$, positivity of this $2\times2$ matrix
and preservation of the unselected marginal give
\begin{equation}
\begin{gathered}
|c_y|\leq\sqrt{pq\,a_yb_y},\\
 c=\sum_y c_y,
 \\
 |c|\leq\sqrt{pq}\sum_y\sqrt{a_yb_y}.
\end{gathered}
\label{eq:chvif-decoder-bhattacharyya}
\end{equation}
This argument retains arbitrary off-diagonal marker blocks and arbitrary
correlations with inaccessible degrees of freedom.

Put $a=a_1$ and $b=b_0$, the two conditional error probabilities. The
readability hypothesis gives $0\leq a,b\leq\varepsilon$, and the last
sum in Eq.~\eqref{eq:chvif-decoder-bhattacharyya} is
\begin{equation}
 \sqrt{(1-a)b}+\sqrt{a(1-b)}.
\end{equation}
Write $a=\sin^2\alpha$, $b=\sin^2\beta$, with
$0\leq\alpha,\beta\leq\arcsin\sqrt{\varepsilon}<\pi/4$. This sum is
$\sin(\alpha+\beta)$, hence it is at most
$2\sqrt{\varepsilon(1-\varepsilon)}$. Multiplying by $2\sqrt{pq}$ gives
Eq.~\eqref{eq:chvif-readable-visibility-bound}, including
$\varepsilon=0$ by continuity or direct substitution.

The weaker bound also follows immediately from weighted discrimination:
the stated decoder has prior-weighted success at least $1-\varepsilon$,
so $D_D=\|p\tau_D^0-q\tau_D^1\|_1\geq1-2\varepsilon$.
The same-state inequality $D_D^2+V^2\leq1$ then gives
$V\leq2\sqrt{\varepsilon(1-\varepsilon)}$. The
Bhattacharyya derivation uses both separate error constraints and therefore
retains the stronger prior-dependent factor. No optimal decoder needs to
be physically available: the actual allowed decoder in
Eq.~\eqref{eq:chvif-worst-label-readability} is enough.

\paragraph{Tight qubit example.}
For any $0\leq\varepsilon<1/2$, take
\begin{align}
 |m_0\rangle&=\sqrt{1-\varepsilon}|0\rangle
                      +\sqrt{\varepsilon}|1\rangle,\\
 |m_1\rangle&=\sqrt{\varepsilon}|0\rangle
                      +\sqrt{1-\varepsilon}|1\rangle,\\
 |\Psi\rangle&=\sqrt p\,|0\rangle|m_0\rangle
                 +e^{i\phi}\sqrt q\,|1\rangle|m_1\rangle.
 \label{eq:chvif-tight-record-qubits}
\end{align}
The computational-basis decoder succeeds with probability
$1-\varepsilon$ on each label, and
$\langle m_1|m_0\rangle=2\sqrt{\varepsilon(1-\varepsilon)}$.
Its visibility saturates the first bound in
Eq.~\eqref{eq:chvif-readable-visibility-bound} for every $p,q$.
For $\varepsilon=0$ the marker is perfectly readable and the route marginal
has zero visibility, while the pure joint state still retains its coherent
superposition. A later joint reversal can recover interference while
changing the record. The theorem bounds the specified unsorted fringe;
it does not identify loss of its contrast with deletion of global coherence.

\subsection{Arrival, retention, and the time of reading}
\label{sec:chvif-record-arrival}

An $(\varepsilon,W_D)$ readable record for domain $D$ requires the
worst-label condition at every allowed reading time $t\in W_D$, with
the channel, allowed decoder, and time origin declared. The theorem then
bounds the contemporaneous visibility at every such time for which the
same route/marker description applies. The statement concerns one read
at a chosen time. A sequence of reads requires the actual reading
instruments and their disturbance; the interval condition alone supplies
no nondestructive repeated-read guarantee.

There is a precise causal pre-arrival restriction. Suppose a newly encoded
label has no label-dependent receiver state in its causal past and the
declared channel has not delivered any label-dependent degree of freedom
by time $t$. Then
\begin{equation}
 \tau_D^0(t)=\tau_D^1(t)=\tau_D(t).
 \label{eq:chvif-before-arrival-states}
\end{equation}
Every binary decoder obeys $r_0+r_1=1$ on this common state, so
\begin{equation}
 \min_j r_j(t)\leq\tfrac12.
 \label{eq:chvif-before-arrival-no-record}
\end{equation}
Thus it cannot satisfy the declared $\varepsilon<1/2$ readability before
arrival. This conclusion needs the explicit identical-state premise,
provided physically by the causal/no-signalling access model. A receiver
that already holds a correlated copy of the label is a different resource
configuration. Unequal priors can improve an average guess before arrival,
but cannot pass the worst-label condition.

The storage window and receiver window are consequently different objects.
Let the end of writing be $t=0$, and suppose the stored encoding is
guaranteed at the memory for emission times $t_e\in[0,T_M]$. For a fixed
latency $\tau$ and a separately budgeted readout time $\tau_{\rm dec}$,
the corresponding idealized completion times are
\begin{equation}
\begin{gathered}
t_D=t_e+\tau+\tau_{\rm dec},\\
 W_D=[\tau+\tau_{\rm dec},\,
                         T_M+\tau+\tau_{\rm dec}].
\end{gathered}
\label{eq:chvif-memory-receiver-windows}
\end{equation}
This translated window assumes that the access channel actually samples
the memory at $t_e$ and that its transport, intervening storage, and decoder
errors have been paid. A protocol needing the memory throughout a finite
read instead requires that entire memory-side interval inside its retention
window. Variable travel times use the declared retarded emission map
$t_e(t_D)$ and its admissible interval. Starting a receiver guarantee at
the end of writing would incorrectly include times before delivery unless
delivery was already available. Choosing the domain time origin at
arrival is valid when that convention is stated.

\subsection{Composable error budgets and retained failed preparations}
\label{sec:chvif-record-composition}

Worst-label budgets compose without independent errors. One convenient
quantum formulation includes all classical outcome flags and relevant
environmental references in the joint state. Let the ideal decoder have
error at most $\varepsilon_0$ on each label. Suppose a hybrid comparison
of the actual preparation, storage, transport, and readout gives
trace-distance errors $\delta_k$ on the relevant joint conditional states
at stage $k$. Each stage bound must hold for the hybrid inputs actually
used, including retained correlations; an entangled reference is included
in the comparison when present. Equivalently, a uniform channel bound
can be supplied on that declared joint input class.

Contractivity under the subsequent trace-preserving channels and the
triangle inequality give a final joint-state distance at most
$\sum_k\delta_k$. A fixed POVM outcome probability changes by no more
than that distance. Consequently
\begin{equation}
\begin{gathered}
r_j^{\rm actual}(t)\geq
       1-\varepsilon_0-\sum_k\delta_k,
 \\
 \varepsilon_{\rm tot}=\varepsilon_0+\sum_k\delta_k.
\end{gathered}
\label{eq:chvif-additive-record-budget}
\end{equation}
The readability and visibility theorem applies when the declared total is
below $1/2$. Isolated marginal error estimates cannot be substituted for
these joint-state bounds when unaccounted correlations change the channel.
The hybrid channels and flags also preserve any feedback used by the
actual instrument.

A classical alternative is an explicit sequence of bad-event flags whose
conditional probability, on every reached prior history and each label,
is at most $\delta_k$. If the final decoder error on the all-good branch
is at most $\varepsilon_0$, the union bound gives the same additive
budget. This argument likewise needs no independence. It does need the
stated conditional bounds; adding unrelated average failure rates would
not establish a worst-label guarantee.

Preparation failure belongs in the full instrument. Suppose label $j$ has
failure probability $f_j\leq f$, and the successful preparation has
worst-label reading error at most $\varepsilon_{\rm succ}$. If an explicit
failure or abstention is counted as an unsuccessful attempted read, then
\begin{equation}
 r_j^{\rm attempt}\geq
       (1-f_j)(1-\varepsilon_{\rm succ})
       \geq1-f-\varepsilon_{\rm succ}.
 \label{eq:chvif-preparation-failure-record}
\end{equation}
Success, failure, and abstention flags remain in the normalized all-outcome
state and in the physical resource account. A larger decoder alphabet can
be retained explicitly, or its failure outcomes can be assigned a stated
binary guess without discarding them.

Alternatively a claim may be conditioned on a heralded success. Its
normalization, conditional record states, and conditional visibility must
then all use that same accepted ensemble. In a route-preserving heralded
preparation the accepted priors become
\begin{equation}
\begin{gathered}
p_{\rm acc}=\frac{p(1-f_0)}{p(1-f_0)+q(1-f_1)},\\
 q_{\rm acc}=1-p_{\rm acc},
\end{gathered}
\label{eq:chvif-accepted-record-priors}
\end{equation}
when the denominator is positive. Outcome-dependent filtering can also
change the coherence, so the accepted joint state itself must be used.
A conditional high-visibility ensemble cannot be compared with a readable
record certified only in a different, unselected ensemble. No reset of
failed preparations or deletion of their records is supplied by this
probability accounting.

\subsection{A finite readable-label capacity consequence}
\label{sec:chvif-record-capacity}

For $N$ declared labels, suppose the full accessible record input has
finite Hilbert dimension $d$ and one POVM $\{F_j\}_{j=1}^N$ satisfies
$\operatorname{Tr}(F_j\sigma_j)\geq1-\varepsilon$ for every normalized
conditional state $\sigma_j$. Since $0\leq\sigma_j\leq I$,
\begin{equation}
 N(1-\varepsilon)
    \leq\sum_{j=1}^N\operatorname{Tr}(F_j\sigma_j)
    \leq\sum_{j=1}^N\operatorname{Tr}F_j=d.
 \label{eq:chvif-approx-readable-capacity}
\end{equation}
Thus $d\geq\lceil N(1-\varepsilon)\rceil$. Any correlated side carrier
available to the decoder belongs in that accessible dimension; alternatively
the decoder is pulled back to a POVM on the declared input space. This
bound concerns $N$ physically decodable labels. Merely naming additional
ordered histories does not add distinguishable record states. It assigns
neither an energy per label nor a thermal erasure cost.

